The criterion ``preserve communities of interest'' appears in redistricting law across
at least thirty states, yet courts that invoke it offer no quantitative definition.
This paper establishes a mathematical framework that operationalizes communities of
interest as cosine similarity of tract-level character vectors, yielding edge weight
modifiers for the METIS graph partitioning engine at the core of the \textsc{bisect}
redistricting system.
The framework is modular: any measurable tract-level signal---economic composition,
land use, housing stock, commuting patterns, topographic setting, administrative zone
co-membership, or transit accessibility---maps to a dimension of the character vector.
The resulting edge weight formula $w(u,v) = w_{\mathrm{base}} \times \mathrm{sim}(\mathbf{c}_u, \mathbf{c}_v)$
satisfies symmetry, non-negativity, and graceful degradation to the unweighted baseline
when all tracts are identical.
We ground the framework in the legal doctrine of communities of interest as articulated
in \textit{Reynolds v.\ Sims} (1964), \textit{Shaw v.\ Reno} (1993), and the independent
redistricting commission criteria enacted in California, Colorado, Arizona, and Washington.
The framework is partisan-neutral: no electoral data, racial composition, or incumbency
information enters any component of the weight computation.
Papers M.1--M.7 implement specific instantiations of the framework; M.8 provides a
composite index.
This paper serves as the foundational theoretical reference for the entire M-track.
